% centralizer_block_diagram.tex
% TikZ block diagram showing the layered structure of Paper II's
% central characterisation result:
%
%   per-site C^12 amplitude
%      |
%      v (discrete-Hermitian centralizer of the bipartite tick rule)
%      |
%   su(6)_+ (+) su(6)_- (+) u(1), dim 71
%      = SM_12 (subalgebra) + extras_59 (SM-module)
%      |
%      v (factor-product projection)
%      |
%   SM gauge algebra + Lorentz (dim 18)

\begin{tikzpicture}[
    every node/.style={font=\small},
    box/.style={
        rectangle, draw, rounded corners=2pt,
        align=center, inner sep=4pt, minimum height=2.4em,
    },
    smbox/.style={
        box, fill=blue!8, draw=blue!50!black, thick,
    },
    extbox/.style={
        box, fill=red!7, draw=red!50!black,
    },
    topbox/.style={
        box, fill=gray!10, draw=gray!50!black, thick,
        minimum width=8.5cm,
    },
    centbox/.style={
        box, draw=black!70, thick,
        minimum width=8.5cm,
    },
    armow/.style={
        -{Latex[length=2.4mm,width=2.0mm]}, semithick,
    },
    label/.style={font=\scriptsize, align=center, gray},
]

% Layer 1: per-site amplitude
\node[topbox] (cone) at (0, 4.4) {%
    per-site amplitude \\
    $\mathbb{C}^{12} \;=\; \mathbb{C}^2_\text{chir} \otimes
     \mathbb{C}^2_\text{iso} \otimes \mathbb{C}^3_\text{col}$
};

% Layer 2: the 71-dim centralizer, decomposed
\node[smbox, minimum width=3.5cm] (sm) at (-2.3, 2.0) {%
    SM subalgebra \\
    $\mathfrak{su}(3)_c \oplus \mathfrak{su}(2)_W \oplus \langle J_1 \rangle$ \\
    \scriptsize dim 12
};
\node[extbox, minimum width=4.6cm] (ext) at (2.45, 2.0) {%
    extras (SM-module, not Lie ideal) \\
    $48 \times (\mathbf{8}, \mathbf{3})_\text{leptoquark}$ \\
    $\;+\; 11 \times \sigma_x\text{-shadow SM}$ \\
    \scriptsize dim 59
};
\node[centbox, fit={(sm) (ext)}, inner sep=8pt, fill=none] (cent) {};
\node[anchor=south, font=\small] at (cent.north) {%
    centralizer \;
    $\mathfrak{su}(6)_+ \oplus \mathfrak{su}(6)_- \oplus \mathfrak{u}(1)$,\quad
    dim 71
};

% Arrow from C^12 to centralizer (with label)
\draw[armow] (cone) -- node[right, label] {%
    Hermitian operators commuting \\
    with $\sigma_x \otimes I_2 \otimes I_3$
} (cent.north -| cone);

% Layer 3: the SM gauge algebra + Lorentz
\node[smbox, minimum width=8.5cm] (eq137) at (0, -1.3) {%
    Eq.~(137) RHS:\quad
    $SO(3,1) \times SU(3) \times SU(2) \times U(1)$,\quad
    dim 18
};

% Arrow from SM subalgebra to Eq. 137 (factor-product projection)
\draw[armow] (sm.south) -- node[left, label, xshift=-2pt] {%
    factor-product \\ projection
} (eq137.north -| sm);

% Arrow from extras showing they're projected to zero
\draw[armow, red!60!black, dashed] (ext.south) -- node[right, label, xshift=2pt] {%
    projected \\ to zero
} ($(eq137.north -| ext) + (0, 0)$);

% Caption fragment showing $L \leftrightarrow R$ bipartite parity
\node[label, anchor=west] at (cent.east) [xshift=4mm, align=left] {%
    bipartite parity \\ $\sigma_x \otimes I_2 \otimes I_3$ \\
    swaps \\ $\mathfrak{su}(6)_+ \leftrightarrow \mathfrak{su}(6)_-$
};

\end{tikzpicture}

\caption[Block diagram: lattice centralizer and the factor-product
projection to the SM gauge algebra.]{%
    The framework's central result.  The per-site amplitude
    $\mathbb{C}^{12}$ (top) admits a discrete-Hermitian
    centralizer (the Hermitian operators commuting with the
    bipartite tick rule's chirality factor $\sigma_x \otimes I_2
    \otimes I_3$) that is structurally $\mathfrak{su}(6)_+ \oplus
    \mathfrak{su}(6)_- \oplus \mathfrak{u}(1)$ of dim 71 (middle),
    with the bipartite parity acting as the $L \leftrightarrow R$
    exchange of the two $\mathfrak{su}(6)$ factors.  This algebra
    decomposes as a non-normal SM gauge subalgebra (left, blue,
    dim 12) plus a 59-dimensional SM-module of ``extras''
    transforming as $(\mathbf{8}, \mathbf{3})$ leptoquark-flavoured
    plus chirality-$\sigma_x$ shadow SM under the SM adjoint
    action.  The factor-product projection (bottom arrow) sends
    each Hermitian generator to its action on a single tensor
    factor of $\mathbb{C}^{12}$ (chirality, isospin, or colour)
    and the other factors to the identity; under this projection
    the SM subalgebra maps onto Eq.~(137)'s right-hand side
    $SO(3,1) \times SU(3) \times SU(2) \times U(1)$ (dim 18) while
    the 59 extras project to zero.  Eq.~(137)'s equality $=$
    is recovered when the factor-product projection is imposed
    as an additional constraint beyond commutativity with the
    bipartite tick rule (sufficiency established; necessity of
    factor-product invariance as the unique additional constraint
    is left open).}
\label{fig:centralizer_block_diagram}
